\section{Introduction}
\label{sec:intro}

Over the last few years, large language models (LLMs) have changed the way software systems encode and operationalize knowledge~\cite{brown2020gpt3,bommasani2021foundation}. Conventional systems concentrate their intellectual value in executable code, where behavior is implemented explicitly. AI agents, by contrast, are increasingly governed by \emph{system prompts}: natural-language specifications of behavior, domain expertise, and safety constraints that act as the primary control interface and externalize system logic into interpretable text~\cite{yao2023react,schick2023toolformer,wang2024llmagents}. A considerable share of the logic that used to reside in compiled artifacts therefore now lives in natural-language specifications. This displacement creates new challenges for intellectual property protection, and their practical relevance grows with the \emph{AI agent as a service} market~\cite{marketsandmarkets2024}.

The problem is most acute in multi-tenant cloud environments, where deployment and data ownership belong to different administrative domains. AI agents are frequently required to execute inside the customer's own cloud account in order to satisfy data residency constraints and regulatory frameworks such as the General Data Protection Regulation (GDPR), the Health Insurance Portability and Accountability Act (HIPAA), and Service Organization Control 2 (SOC~2)~\cite{gdpr2016}. At the same time, the system prompt must be treated as proprietary knowledge, because it encodes the differentiated capabilities on which the commercial value of the service rests. The two requirements conflict: running the agent in a customer-controlled environment opens an exposure surface in which prompt configurations become readable by privileged actors, which defeats their confidentiality~\cite{aws_shared_responsibility,agenticsecurity2025}.

The conflict is practical rather than hypothetical. Managed agent platforms such as Amazon Bedrock Agents, Azure AI Agent Service, and Vertex AI Agents keep agent configurations in account-level resources that administrative principals can usually read. Prompt extraction has been demonstrated without any privileged access~\cite{hui2024pleak,multiturn2024}, and prompt injection~\cite{promptinjection2026,owasp2025}, ranked first in the Open Worldwide Application Security Project (OWASP) LLM Top-10, enlarges the attack surface further by enabling indirect manipulation of model behavior.

Existing defenses cover this problem only partially and, as a rule, in isolation. Cryptographic mechanisms protect prompts in transit or under controlled access~\cite{vanWyk2023protect}. Software defenses address individual vectors through filtering, obfuscation, or leakage detection~\cite{proxyprompt2025,sysvec2025,encryptedprompt2025,promptguard2026,promptkeeper2024}. Hardware-based approaches rely on Trusted Execution Environments (TEEs), which offer strong isolation but require dedicated hardware, add measurable overhead~\cite{teeperf2025}, and provide no support for forensic attribution~\cite{petridish2024,omega2025}. Surveys of agentic security~\cite{agenticsecurity2025} document the breadth of the threat landscape, but they report no deployed architecture that protects prompts systematically across administrative domains.

We therefore frame prompt protection in multi-tenant LLM systems as a cross-domain trust problem rather than as the design of one more isolated security mechanism~\cite{saltzer1975protection,rose2020zero,agenticsecurity2025}. This framing calls for reasoning about how trust is split between vendor and customer environments, how adversaries exploit the resulting trust boundary, and how several defensive mechanisms can be composed so that the security guarantees survive the compromise of any single one.

This article presents a defense-in-depth architecture that composes standard cloud services into a four-layer trust chain: cryptographic prompt delivery (Layer~1), OWASP-aligned input validation (Layer~2), output filtering with steganographic watermarking (Layer~3), and cross-account monitoring (Layer~4). The composition provides prompt confidentiality, output integrity, and forensic traceability while preserving deployment flexibility in multi-tenant settings. Since it depends on no specialized hardware and distributes the defenses across independent layers, the architecture is portable across cloud providers and tolerant to the failure of individual layers.

A layered design, however, raises two questions that a single deployment cannot answer. Layers~2 and~3 interact with the behavior of the underlying LLM, so their effectiveness may vary across model families; and the value of composing several layers rests on the assumption that they are complementary, which has to be shown empirically. We formulate these questions as follows.

\begin{itemize}
\item \textbf{RQ1 (model robustness).} How robust is the proposed defense-in-depth architecture across different LLM families?
\item \textbf{RQ2 (layer contribution).} How much does each defense layer contribute to the overall security effectiveness of the architecture?
\end{itemize}

This article is an extended and revised version of our work presented at JCC-BD\&ET 2026~\cite{petrocelli2026jcc}, which introduced the threat model, the four-layer architecture, the serverless reference implementation, and its validation through controlled penetration testing and a live two-account deployment. The present version extends that work with the following contributions: (i)~a systematic multi-model evaluation that applies an identical experimental protocol to nine LLMs from six model families, both proprietary and open-weight, in order to answer RQ1; (ii)~a defense-layer ablation study that quantifies the individual and complementary contribution of the defense layers, including the cases in which an attack that passes Layer~2 is intercepted by Layer~3, in order to answer RQ2; and (iii)~an expanded experimental methodology, a more detailed account of the reference implementation and of its reproducibility environment, and a discussion informed by the new evaluations. The sections shared with the conference version have been revised and rewritten.

The evaluation covers three complementary dimensions: (i)~security effectiveness, understood as resistance to prompt extraction and adversarial manipulation; (ii)~performance overhead, measured as the latency added by the defense layers; and (iii)~practical deployability, including portability across providers and feasibility under realistic operational constraints.

The remainder of this article is organized as follows. Section~\ref{sec:related} reviews related work. Section~\ref{sec:threats} formalizes the threat model and the security requirements. Section~\ref{sec:arch} describes the architecture, and Section~\ref{sec:implementation} its reference implementation and reproducibility environment. Section~\ref{sec:methodology} presents the experimental methodology, Section~\ref{sec:security-eval} the security evaluation, and Section~\ref{sec:perf-cost} the performance and cost evaluation. Section~\ref{sec:discussion} discusses the results and their limitations, and Section~\ref{sec:conclusions} concludes the article.
